% ##################################################################################################################
\section{Padang}
\label{sec:padang}
\hfill \textbf{Author:} Gregor Lämmel

\editdone{This text has undergone the professional edit. Please no grammatical changes anymore! They are most-probably wrong.}

% ##################################################################################################################
The Padang scenario demonstrated the \gls{matsim} application to large-scale evacuation problems. The scenario was created as part of the third party funded project, "Last-Mile". \citet{TaubenboeckEtAl2012ConcludingLastMilePaperNatHazards} gives a comprehensive overview.
Padang is located on the west coast of Sumatra Island, Indonesia. In 2014, the city had a population of about 1\,000\,000 people. 
Because of its problematic location on the coast in a so-called "seismically locked" area~\citep{McCloskey2010Padang2009Earthquake}, Padang is prone to earthquakes and subsequent tsunamis. In the "Last-Mile" project, a realistic tsunami scenario, triggered by an earthquake about 300\,km off the coast, was identified~\citep{GosebergSchlurmann2009HazardMappingPadang}. The assumed tsunami would have left about 30\,minutes for the evacuation. The flooding would reach as far as 3\,kilometers inland, thus threatening up to 330\,000 lives. \citet{Laemmel_PhDThesis_2011} developed a \gls{matsim} scenario representing the city with its affected population. One unusual aspect of the Padang situation was the expected universal evacuation by foot; simulating pedestrians with \gls{matsim} was a novelty when this project was undertaken. The standard simulation model (see, e.g.,\,Chapter~\ref{sec:trafficflowmodel}) was thus adapted to deal with pedestrians. 
Details were discussed in \citet{LaemmelKluepfelNagel2009EvacPadangAtBookTimmermanns}. Another important variation, contrary to most standard transport scenarios, was that network links would flood once the tsunami reached them. Thus, accessibility---flooded or not flooded---of the network links was time-dependent, which was modeled by special time-dependent networks \citep{LaemmelEtAl_TransResC_2010}. In the time-dependent network concept, link attributes---like \emph{freespeed}---could be changed, while the simulation ran, by precomputed network change events. For the Padang scenario, the network change events were extracted from microscopic flooding simulation data.

% ##################################################################################################################
%
\createfigure%
{Evacuation time analysis for downtown Padang. Numbers showing average evacuation time in minutes, which are also indicated by the colors green, yellow, red.}%
{Evacuation time analysis for downtown Padang. Numbers showing average evacuation time in minutes, which are also indicated by the colors green, yellow, red.}%
{\label{chap:using:padang}}%
{\includegraphics[width=0.6\textwidth, angle=0]{using/figures/dwntwnpdg}}%
{}

Key Padang scenario facts:
\begin{compactitem}
\item The network consisted of about 6\,000 nodes and 17\,000 links.
\item Synthetic populations for morning, afternoon, and night were created, containing up to 330\,000 agents.
\item The flooding was modeled by a set of 109\,network change events (one per minute), affecting 7\,609 links.
\item A set of 42\,shelter buildings, which could be used for vertical evacuation, was also part of the scenario.
\end{compactitem}
Based on the Padang scenario, various evacuation strategies were investigated:
\begin{compactitem}
\item A seemingly obvious evacuation strategy was the shortest path solution, where everyone is on the shortest path. This solution, however, ignored possible congestion and led to unfeasible results:
\item The Nash equilibrium approach was better, where everyone tried to find an optimal evacuation route through iterative learning \citep{LaemmelKluepfelNagel2009EvacPadangAtBookTimmermanns};
\item While the Nash equilibrium reduced individual evacuation time, total evacuation time might not have been minimized. The marginal social cost-based simulation approach tried to minimize total evacuation time \citep{LaemmelFloetteroed_MertschingEtAl_2009,DresslerFloetteroedLaemmelNagelSkutella2010OptimalEvacuationLargeScaleScenarios};
\item These three basic evacuation approaches were investigated together with flooding \citep{LaemmelEtAl_TransResC_2010,Laemmel_PhDThesis_2011};
\item Further, an evacuation strategy to reduce exposure risk was developed by \citep{LaemmelKluepfelNagel2010PEDRiskPrinted};
\item And finally, \citet{FloetteroedLaemmel2010ICECShelterEvac} proposed a method to integrate shelter buildings, which are evacuation sinks (i.e.,\,safe places) 
% "\Karen{ What?? sink? Does he mean 'shed', or some kind of building? Or do the buildings have sanitary facilities? Need help with this word.}" 
with limited capacity, into the simulation.  
\end{compactitem}

% ##################################################################################################################
